\documentclass[11pt]{article}
\usepackage[a4paper,margin=1in]{geometry}
\usepackage{subfiles}
\usepackage{amsmath,amssymb,bm}
\usepackage{hyperref}
\usepackage{booktabs}
\usepackage{array}
\usepackage{mathtools}

\newcommand{\rs}{r_{\mathrm{s}}}

\title{Appendix 13 (to main theory): Covariant Hysteresis Sector: Augmented Action, Stability, and Quasi-Static Reduction}
\author{}
\date{}

\begin{document}
\let\phi\varphi
\maketitle

\section*{Purpose, scope, and reader contract}
This appendix accompanies the main theory. It serves two purposes:
\begin{enumerate}
\item \textbf{Baseline map:} a compact dependency map of the ISPG ingredients used downstream;
\item \textbf{Hysteresis phenomenology:} a specification of effective memory effects (captured by the field $\chi$ or, at effective level, by the transported hysteretic response derived in the companion MOND paper~\cite{Lotishvili2026MOND}) that reproduce galaxy-halo phenomenology without new fundamental particle species.  In the covariant template below, $\chi$ is an effective collective memory field, not a new fundamental degree of freedom.
\end{enumerate}

\textbf{Relation to the companion MOND paper~\cite{Lotishvili2026MOND}.}
the companion MOND paper~\cite{Lotishvili2026MOND} derives the MOND interpolating function $\mu(x)=x/(1+x)$ from the bound-structure Bernoulli$+$vortex closure in rotating galaxies, with the scalar field equation $\Box\varphi=S$ on a weak-field rotating background fixing the relaxation time scale and transport rate. The present Appendix~13 provides an \emph{effective closure} perspective on the hysteretic part of the same physics: the field $\chi$ captures the transported memory response that the companion MOND paper~\cite{Lotishvili2026MOND} derives from first principles via Fourier decomposition and Bessel eigenvalue analysis. The two appendices are \emph{complementary}: the companion MOND paper~\cite{Lotishvili2026MOND} gives the formal bound-structure derivation; Appendix~13 gives the phenomenological implementation and strengthening routes for future work.
In other words, the $\chi$ appearing here is not a second unrelated MOND coefficient: it is the dynamical coarse-grained field-level representative of the same transported response that the companion MOND paper~\cite{Lotishvili2026MOND} writes locally as the vortex-source multiplier.  The precise mathematical identification of $\chi$ with the companion paper's local multiplier---the explicit mapping between Eq.~\eqref{eq:chi_KG} and the vortex-closure equation at the bound-structure level---is the content of the Route~B coarse-graining programme (Sec.~\ref{sec:ledger_app13}) and is listed there as an open item; the statement here is that both objects describe the same physical quantity, not that the mapping has been carried out in closed form.

\textbf{Scope boundary.}
The $\chi$ evolution laws were originally introduced as
an effective closure reproducing the results of the companion MOND paper~\cite{Lotishvili2026MOND}.
In Sec.~\ref{sec:routeA_real} of this appendix,
\textbf{Route~A is treated as a covariant effective template}
for an effective collective field with standalone covariant
dynamics in the quasi-static bound-structure sector.  The
standalone $\chi$ equation is ghost-free, gradient-stable,
and normally hyperbolic when its source is treated as an
external quasi-static functional of the background scalar.
The fully varied backreacting $(\varphi,\chi)$ principal-symbol
analysis remains open.  The main strengthening route is:
\begin{itemize}
\item \textbf{Route B (coarse--graining/EFT):} obtain $\chi$ as an emergent slow variable from the nonlinear dynamics of $\phi$ around bound structures, connecting to the rotational bound-structure analysis of the companion MOND paper~\cite{Lotishvili2026MOND}.
\end{itemize}

\section{Baseline: bi--conformal metric and covariant action}\label{sec:base}
\subsection{Bi--conformal metric ansatz}
The constructive requirement is that a single scalar $\phi$ controls local clock rates and local spatial scales reciprocally. This is implemented by the bi--conformal ansatz
\begin{equation}
ds^2 = -e^{\phi}\,c^2 dt^2 + e^{-\phi}\left(dx^2+dy^2+dz^2\right).
\label{eq:biconf_app13}
\end{equation}

\subsection{Covariant action and minimal Horndeski placement}
The baseline action is of scalar--tensor form with minimal coupling of matter to $g_{\mu\nu}$:
\begin{equation}
S=\frac{1}{16\pi G}\int d^4x\sqrt{-g}\left[R+\frac12 g^{\mu\nu}\partial_\mu\phi\,\partial_\nu\phi\right]+S_m[g_{\mu\nu},\psi].
\label{eq:action_app13}
\end{equation}
A compact Horndeski mapping (minimally coupled single--field subclass) is:
\begin{equation}
G_2 = -\frac{M_{\rm Pl}^2}{2}X,\qquad
G_3=0,\qquad
G_4=\frac{M_{\rm Pl}^2}{2},\qquad
G_5=0,\qquad
X\equiv -\frac12 \partial_\mu\phi\,\partial^\mu\phi.
\label{eq:horndeski_app13}
\end{equation}
(Discussion of the sign choice and stability conditions is given in the main theory and is not reopened here.)

\section{Weak--field limit and the Shapiro discriminator}\label{sec:ppn_app13}
Expanding \eqref{eq:biconf_app13} for small $\phi$ yields
\begin{equation}
g_{tt}\simeq -(1+\phi),\qquad g_{ij}\simeq (1-\phi)\delta_{ij},
\end{equation}
so that $\Phi=\Psi=(c^2/2)\phi$ and $\gamma_{\rm PPN}=1$ at leading order.

The main theory derives the 1PN Shapiro delay via an effective refractive index $n(r)=e^{-\phi}$ (identical to GR at 1PN, as $\gamma=1$) and identifies a structurally distinct 2PN signature from the exponential $n=e^{-\phi}$; the refractive sub-piece of the 2PN delay is derived exactly and the closed ISPG--GR difference $\Delta B = \pi/4$ is established at the regularization-independent level (the main theory Shapiro section~\cite{Lotishvili2026} and Appendix~6).  Standalone absolute 2PN coefficients remain convention-dependent finite-endpoint quantities, but the differential discriminator used by the theory is closed. For later use we record the dependency structure:
\begin{equation}
n(r)=e^{-\phi}\simeq 1+\frac{\rs}{r},\qquad
\Delta t_{2\rm PN}^{\rm ISPG}\propto \frac{G^2M^2}{c^5 b}.
\end{equation}

\section{Linear cosmology versus bound structures: regime separation}\label{sec:regime}
A central design goal of the main theory is that \emph{linear} cosmology can remain GR-identical in the same-matter-content metric-sector limit while \emph{nonlinear} bound structures exhibit new behaviour. The separation is operationally expressed by a scale hierarchy:
\begin{equation}
k \ll \frac{aH}{c}\quad (\text{super-/horizon scales, linear cosmology}),\qquad
k \gg \frac{aH}{c}\ \text{and}\ \partial_t \ll c\nabla\quad (\text{quasi-static, bound structures}).
\label{eq:scaleseparation}
\end{equation}

\textbf{Cosmic-time compatibility.}
In the cosmic-time frame tied to matter clocks (comoving observers), $t=\tau$ along comoving worldlines, so $g_{tt}=-1$ in standard FLRW form. The bi--conformal structure \eqref{eq:biconf_app13} then requires $e^{\phi_0}=1$ and hence $\phi_0=0$ at the homogeneous level. This fixes the \emph{background} without eliminating \emph{inhomogeneous} responses sourced by matter in the quasi-static sector.

\section{Matter-sourced scalar response in the quasi-static sector}\label{sec:deltaphi}
The representative matter-sourced perturbation equation used in the main theory (metric-decoupled at linear order under $\alpha_i=0$, but not a free spectator; it is driven by $\delta T_k$ and inherits the phantom-sign caveat of the background action, cf.\ App.~11) is
\begin{equation}
\delta\ddot{\phi}_k + 3H\,\delta\dot{\phi}_k + \left(\frac{c^2k^2}{a^2}\right)\delta\phi_k = \frac{8\pi G}{c^2}\,\delta T_k.
\label{eq:deltaphi_app13}
\end{equation}
The RHS coefficient $8\pi G/c^2$ follows from the FLRW expansion of the sourced scalar equation $\Box\phi=-(8\pi G/c^4)T$ of App.~1: $\Box\phi=-c^{-2}(\ddot\phi+3H\dot\phi)+a^{-2}\nabla^2\phi$, multiplied by $-c^2$, gives the form above and matches App.~9 Eq.~\eqref{eq:deltaphi_app9}.
Equation~\eqref{eq:deltaphi_app13} makes explicit that localized inhomogeneities can excite $\delta\phi$ through $\delta T_k$ even when the homogeneous background is fixed to $\phi_0=0$ in cosmic time. The MOND phenomenology of the main theory is associated with the quasi-static, weak-field, nonrelativistic limit around bound structures, not with the linear FLRW sector.

\section{The hysteresis field $\chi$: closure, selection principle, and strengthening routes}\label{sec:chi}
\subsection{Intent and modelling posture}
We introduce $\chi$ as the effective field capturing medium--memory (hysteresis) effects. In this appendix it serves as the field-level effective closure of the transported response discussed in the companion MOND paper~\cite{Lotishvili2026MOND}. Its purpose is to (i) encode memory effects with minimal extra structure, (ii) produce falsifiable halo-scale consequences, and (iii) admit a controlled coarse-graining derivation from the underlying bound-structure dynamics via Route~B.

\subsection{Phenomenological relaxation ansatz}
As motivation, we first record the relaxation law toward a local equilibrium functional $\chi_{\rm eq}$:
\begin{equation}
\dot{\chi} = -\tau^{-1}\big(\chi-\chi_{\rm eq}\big),
\label{eq:chi1}
\end{equation}
with $\tau$ to be determined.  (This ansatz is \emph{recovered}
as an overdamped, coarse-grained limit of the covariant template
in Sec.~\ref{sec:qsreduction}; the free massive KG sector by itself
oscillates rather than relaxing.)

\subsection{Closure choice for relaxation}
We take the leading local equilibrium to be proportional to the quasi-static gradient strength,
\begin{equation}
\chi_{\rm eq}=\beta\,|\nabla\phi|,\qquad \beta>0,
\label{eq:chieq}
\end{equation}
and we adopt the leading local relaxation-rate law
\begin{equation}
\tau^{-1} = \gamma\,|\nabla\phi|\qquad \text{(Variant A)}.
\label{eq:chi2}
\end{equation}

\paragraph{Working assumptions (explicit).}
In the hysteresis sector we restrict attention to the quasi-static, weak-field regime around bound structures, characterized by $\partial_t \ll c\nabla$ and (for Fourier modes) $k\gg aH/c$. The closure is required to be \emph{local}: it depends only on local scalars built from the fields and the metric in this regime. Concretely, at leading order we allow dependence on
\begin{equation}
|\nabla\phi|,\qquad
Y_{\rm qs}\equiv \sqrt{\nabla_\mu\phi\,\nabla^\mu\phi},
\end{equation}
where $Y_{\rm qs}$ is defined only in the spacelike-gradient,
quasi-static bound-structure regime.  We exclude explicit dependence
on coordinate radius $r$, galaxy-specific input scales, or ad hoc
profile functions. The ``no new macroscopic scales'' requirement means
that $\tau$ and $\chi_{\rm eq}$ must be expressible in terms of these
local invariants and constants already present in the baseline sector,
without introducing an additional halo-tuned length or time scale.

\paragraph{Minimal (leading-order) character of Variant A.}
Under locality, monotonicity (stronger gradients relax faster), and the absence of additional macroscopic scales, the choice $\tau^{-1}\propto|\nabla\phi|$ can be viewed as the leading-order term in a local gradient-expansion for the relaxation rate. Alternative closures correspond to higher-order corrections (involving additional invariants or suppressed operators) or to genuinely new microphysical scales; in either case they predict quantitatively distinct halo profiles and lensing--dynamics relations and are therefore observationally discriminable.

\subsection{Static limit $\Rightarrow$ isothermal scaling (consistency check)}
In the static weak-field limit of isolated sources one has $|\nabla\phi|\propto 1/r^2$. Under Variant~A, the closure drives $\chi\propto |\nabla\phi|$ and therefore $\chi\propto 1/r^2$ (and similarly $\rho_\chi\propto 1/r^2$ if $\rho_\chi$ is taken proportional to $\chi$).
\textbf{Intended reading:} this is a \emph{consistency check} for the chosen closure, not a first-principles derivation. The decisive content is that other relaxation laws predict different halo profiles, providing observational discriminants.

\subsection{Rotation curves and lensing fingerprints}
Because $\chi$ is tied to gradients of $\phi$, galaxy dynamics and lensing are linked. Different closure choices $\tau(|\nabla\phi|)$ predict different outer-halo slopes and transition radii, and therefore distinct relations between rotation curves and shear profiles. This is one of the primary falsification channels.

\subsection{Route A sketch: a concrete covariant template}\label{sec:routeA}
To make Route~A concrete, we record a covariant effective template in
which a memory field is driven toward a gradient-built equilibrium.
Throughout this Route~A section, the source functional is treated as
an external quasi-static functional of the background $\phi$; the fully
varied two-field principal-symbol analysis is an open task.  Define
the quasi-static invariant
\begin{equation}
Y_{\rm qs} \equiv \sqrt{\nabla_\mu\phi\,\nabla^\mu\phi}\quad
\text{(spacelike-gradient dominated in the quasi-static sector)}.
\end{equation}
A minimal augmented-action ansatz can be written schematically as
\begin{equation}
S[g,\phi,\chi]=S_{\rm ISPG}[g,\phi] + \int d^4x\sqrt{-g}\left[-\frac{1}{2}(\nabla\chi)^2 - U(\chi) - \frac{\lambda}{2}\big(\chi-\beta Y_{\rm qs}\big)^2\right],
\label{eq:augaction}
\end{equation}
where $\lambda$ controls the stiffness toward the local equilibrium
$\chi\simeq \beta Y_{\rm qs}$. In the quasi-static sector with
sufficiently large $\lambda$, the field tracks
$\chi_{\rm eq}\approx \beta Y_{\rm qs}$ up to controlled corrections.
We now record the supported one-way effective construction explicitly.

\section{Route~A realization: field equation and stability}\label{sec:routeA_real}

\subsection{Minimal action ($U = 0$)}\label{sec:minaction}
Following the parameter-economy principle, we set $U(\chi) = 0$
in~\eqref{eq:augaction}, retaining only the kinetic and
stiffness terms:
\begin{equation}
S_\chi = \int d^4x\,\sqrt{-g}\left[
  -\frac{1}{2}\,\nabla_\mu\chi\,\nabla^\mu\chi
  - \frac{\lambda}{2}\bigl(\chi - \beta\,Y_{\rm qs}\bigr)^2
\right],
\label{eq:Schi}
\end{equation}
with
\begin{equation}
Y_{\rm qs} \equiv \sqrt{\nabla_\mu\varphi\,\nabla^\mu\varphi}
\end{equation}
in the spacelike-gradient, quasi-static bound-structure sector,
with $\lambda > 0$ carrying dimensions of
$[\mathrm{length}^{-2}]$ (equivalently mass$^{2}$ in
natural units) and $\beta > 0$.
Here $\chi$ is an effective collective field with standalone covariant
dynamics in the quasi-static bound-structure sector. In that limit it
should be read as the field-level effective realization of the local
transported-response multiplier used in the companion MOND
paper~\cite{Lotishvili2026MOND}, while Route~B remains the microscopic
coarse-graining step.
No additional self-interaction potential is introduced;
the only scale controlling $\chi$ dynamics is the
stiffness parameter~$\lambda$.

\paragraph{Microscopic stiffness scale.}
In the classical context of this appendix, $\lambda$ was
originally a closure stiffness.  The quantum
extension~\cite{Lotishvili2026Q} identifies the microscopic tracking
scale with the quasi-normal mode structure of the oscillon:
\begin{equation}
  \lambda_{\mathrm{micro}}
  = \frac{\omega_0^2}{c^2}
  = \left(\frac{m_0 c}{\hbar}\right)^{\!2} e^{\varphi}\,,
  \label{eq:lambda_micro}
\end{equation}
where $\omega_0$ is the fundamental confining frequency
and $m_0$ the oscillon rest mass.  This fixes the
microphysical tracking stiffness, but it does not by itself derive the
macroscopic collective relaxation scale used in galaxy-scale
hysteresis.  That macro-scale matching remains part of the Route~B
coarse-graining programme, and $\beta$ remains a phenomenological
closure normalization until the full bound-structure derivation is
completed.

\paragraph{Domain and regularity note.}
The source functional
$Y_{\rm qs} = \sqrt{\nabla_\mu\varphi\,\nabla^\mu\varphi}$
is real only in the spacelike-gradient or null-gradient regime and is
therefore used here as a quasi-static bound-structure quantity, not as
a homogeneous FLRW/timelike-gradient scalar.  It is also non-smooth at
$Y_{\rm qs} = 0$: the variation
$\delta Y_{\rm qs}/\delta(\partial_\mu\varphi) =
\nabla^\mu\varphi/Y_{\rm qs}$
diverges there.  In the quasi-static sector around
gravitating sources $\nabla\varphi \neq 0$ everywhere
outside the source centre, so the singularity is not
encountered in practice.  Should a fully global treatment
be required, the standard remedy is to replace $Y_{\rm qs}$ with
a regularized form
$Y_{{\rm qs},\epsilon} \equiv
\sqrt{\nabla_\mu\varphi\,\nabla^\mu\varphi
+ \epsilon^2}$ ($\epsilon \to 0^+$), which preserves all
quasi-static results derived below while rendering the action smooth.

\subsection{Field equation for $\chi$}\label{sec:chi_eom}
Varying~\eqref{eq:Schi} with respect to $\chi$ gives
\begin{equation}
\frac{\delta S_\chi}{\delta\chi}
= \Box\chi - \lambda\bigl(\chi - \beta\,Y_{\rm qs}\bigr) = 0\,,
\label{eq:chi_eom}
\end{equation}
which can be written as
\begin{equation}
\boxed{\Box\chi - \lambda\,\chi = -\lambda\,\beta\,Y_{\rm qs}\,.}
\label{eq:chi_KG}
\end{equation}
This is a Klein--Gordon equation with effective
mass $m_\chi^2 = \lambda$ and source $\lambda\beta Y_{\rm qs}$.
The source is built entirely from the scalar field
$\varphi$ of the baseline ISPG sector; no new coupling
to matter is introduced.

\paragraph{Structure of the equation.}
\begin{itemize}
\item In regions where $\varphi$ is slowly varying
  ($Y_{\rm qs} \to 0$), the source vanishes.  The free
  homogeneous KG mode then oscillates with frequency
  $c\sqrt{\lambda}$; relaxation to zero requires
  damping, coarse-graining, or outgoing boundary
  conditions.
\item Near a gravitating source where
  $|\nabla\varphi| \sim r_s/r^2$, the field is driven
  toward $\chi \approx \beta Y_{\rm qs}$, i.e.\ the local
  equilibrium of the phenomenological closure.
\item The $\Box\chi$ term provides the propagation dynamics
  absent in the original relaxation ansatz~\eqref{eq:chi1},
  upgrading it from a first-order ODE to a covariant
  hyperbolic PDE.
\end{itemize}

\subsection{Ghost freedom}\label{sec:ghost_app13}
A ghost (negative-norm excitation) arises when the
kinetic term of a dynamical field carries the wrong
sign, rendering the Hamiltonian unbounded from below.
In~\eqref{eq:Schi} the $\chi$~kinetic term is
\begin{equation}
\mathcal{L}_{\rm kin}^{(\chi)}
= -\tfrac{1}{2}\,\nabla_\mu\chi\,\nabla^\mu\chi\,.
\label{eq:Lkin}
\end{equation}
With the mostly-plus signature convention
$(-,+,+,+)$ used throughout the main theory,
$\nabla_\mu\chi\,\nabla^\mu\chi
= -\dot\chi^2/c^2 + (\partial_i\chi)^2$,
so the Lagrangian density becomes
\begin{equation}
\mathcal{L}_{\rm kin}^{(\chi)}
= +\frac{1}{2c^2}\,\dot\chi^2
  - \frac{1}{2}\,(\partial_i\chi)^2\,.
\label{eq:Lkin_3p1}
\end{equation}
The time-kinetic part $+\dot\chi^2/(2c^2)$ enters
with a \emph{positive} sign.  The conjugate momentum is
\begin{equation}
\pi_\chi = \frac{\partial\mathcal{L}}{\partial\dot\chi}
= \frac{\dot\chi}{c^2}\,,
\label{eq:pi_chi}
\end{equation}
and the corresponding Hamiltonian density is
\begin{equation}
\mathcal{H}_\chi
= \frac{c^2}{2}\,\pi_\chi^2
  + \frac{1}{2}\,(\partial_i\chi)^2
  + \frac{\lambda}{2}\,(\chi - \beta Y_{\rm qs})^2\,.
\label{eq:H_chi}
\end{equation}
Every term is non-negative for $\lambda > 0$.
Therefore the $\chi$~sector is \textbf{ghost-free}:
the energy is bounded from below and no negative-norm
states propagate.

\subsection{Gradient stability}\label{sec:gradient}
Gradient instability occurs when spatial perturbations
grow exponentially, signalling that the field equation
is ill-posed as an initial-value problem.
To check this, linearize~\eqref{eq:chi_KG} around a
static background $\bar\chi$:
\begin{equation}
\chi = \bar\chi + \delta\chi\,,\qquad
\delta\chi \propto e^{-i\omega t + i\mathbf{k}\cdot\mathbf{x}}\,.
\label{eq:pert_ansatz}
\end{equation}
Substituting into~\eqref{eq:chi_KG} and keeping only
the homogeneous part (the source $\lambda\beta Y_{\rm qs}$ is
fixed by the background $\varphi$), the dispersion
relation is
\begin{equation}
\boxed{\omega^2 = c^2\,|\mathbf{k}|^2 + \lambda\,c^2\,.}
\label{eq:dispersion_app13}
\end{equation}
For all $\lambda > 0$:
\begin{itemize}
\item $\omega^2 > 0$ for every wavenumber $\mathbf{k}$,
  so no mode grows exponentially.
\item The group velocity
  $v_g = \partial\omega/\partial|\mathbf{k}|
  = c^2|\mathbf{k}|/\omega < c$,
  so $\chi$ perturbations propagate subluminally
  (massive field).
\item The phase velocity $v_p = \omega/|\mathbf{k}|
  = c\sqrt{1 + \lambda/|\mathbf{k}|^2} > c$ in the
  usual massive-field sense, but carries no information
  superluminally.
\end{itemize}
The spatial-gradient coefficient in the
Hamiltonian~\eqref{eq:H_chi} is $+\tfrac{1}{2}$
(positive definite), confirming the absence of
gradient instability at the nonlinear level as well.

\subsection{Quasi-static reduction to the phenomenological
  closure}\label{sec:qsreduction}
We now show that the covariant field
equation~\eqref{eq:chi_KG} reproduces the
phenomenological relaxation law~\eqref{eq:chi1}
in the appropriate limit.

In the FLRW background the d'Alembertian acquires
a Hubble friction term:
\begin{equation}
\Box\chi = -\frac{1}{c^2}\,\ddot\chi
  - \frac{3H}{c^2}\,\dot\chi
  + \frac{1}{a^2}\,\nabla^2\chi\,.
\label{eq:box_flrw}
\end{equation}
The field equation~\eqref{eq:chi_KG} becomes
\begin{equation}
-\frac{1}{c^2}\,\ddot\chi
  - \frac{3H}{c^2}\,\dot\chi
  + \frac{1}{a^2}\,\nabla^2\chi
  + \lambda\,\chi
  = \lambda\,\beta\,Y_{\rm qs}\,.
\label{eq:chi_flrw}
\end{equation}

\paragraph{Overdamped, long-wavelength limit.}
In the quasi-static sector around bound structures,
two conditions hold:
\begin{enumerate}
\item \emph{Overdamped:} the field evolves on
  timescales much longer than $1/\sqrt{\lambda}\,c$,
  so $|\ddot\chi| \ll \lambda c^2\,|\chi|$ and
  the $\ddot\chi$~term can be dropped.
\item \emph{Long-wavelength:} the characteristic
  spatial scale $\ell$ of $\chi$ satisfies
  $1/\ell^2 \ll \lambda$, so
  $|\nabla^2\chi| \ll \lambda\,|\chi|$ and
  the Laplacian can be dropped.
\end{enumerate}
Under these two conditions, Eq.~\eqref{eq:chi_flrw}
reduces to
\begin{equation}
-\frac{3H}{c^2}\,\dot\chi
  = \lambda\bigl(\chi - \beta\,Y_{\rm qs}\bigr)\,,
\end{equation}
which can be rewritten as
\begin{equation}
\boxed{\dot\chi
  = -\frac{\lambda\,c^2}{3H}
    \bigl(\chi - \beta\,Y_{\rm qs}\bigr)\,.}
\label{eq:chi_relaxation}
\end{equation}
This is precisely the phenomenological
relaxation~\eqref{eq:chi1} with the identifications
\begin{equation}
\chi_{\rm eq} = \beta\,Y_{\rm qs}\,,\qquad
\tau^{-1} = \frac{\lambda\,c^2}{3H}\,.
\label{eq:tau_derived}
\end{equation}

\paragraph{Physical content.}
\begin{itemize}
\item Within the overdamped, long-wavelength reduction,
  the effective relaxation timescale
  $\tau = 3H/(\lambda c^2)$ is fixed by the ratio of
  Hubble friction to the stiffness parameter~$\lambda$.
\item For $\lambda \gg H^2/c^2$ the field tracks $\beta Y_{\rm qs}$
  almost instantaneously; for $\lambda \sim H^2/c^2$
  the lag becomes cosmologically significant, producing
  observable hysteresis.
\item The equilibrium value $\chi_{\rm eq} = \beta Y_{\rm qs}
  \approx \beta\,|\nabla\varphi|$ in the quasi-static
  sector, matching the phenomenological
  closure~\eqref{eq:chieq}.
\item The original Variant~A ansatz
  $\tau^{-1} = \gamma\,|\nabla\varphi|$
  (Eq.~\ref{eq:chi2}) was a phenomenological
  gradient-dependent closure originally introduced to
  reproduce the isothermal scaling of the halo regime.
  The present covariant derivation supersedes it: the
  leading-order relaxation rate obtained directly from
  the action~\eqref{eq:Schi} is instead
  Hubble-regulated, $\tau^{-1}=\lambda c^2/(3H)$, and
  is not gradient-dependent to this order.
  Reconciling the two forms therefore requires either
  a subleading gradient-dependent correction to the
  covariant $\lambda$, or the identification of a
  specific regime in which the effective Variant~A
  re-emerges from the covariant theory.  This
  reconciliation is listed as an open item in the
  Route~B programme (Sec.~\ref{sec:ledger_app13}),
  and the outer-halo region $|\nabla\varphi|$ varies
  rapidly provides the natural observational setting
  in which any residual gradient dependence would be
  tested against the Hubble-regulated baseline.
\end{itemize}

\subsection{Hyperbolicity and well-posedness}\label{sec:wellposed}
A PDE system is well-posed as an initial-value problem
if and only if its principal part is hyperbolic.
The principal symbol of the $\chi$ field
equation~\eqref{eq:chi_KG} is determined by its
highest-derivative term:
\begin{equation}
\mathcal{P}(\chi) = \Box\chi + (\text{lower order})
\quad\Longrightarrow\quad
\sigma(\xi) = g^{\mu\nu}\xi_\mu\xi_\nu\,,
\label{eq:principal_app13}
\end{equation}
where $\xi_\mu$ is the covector of the Fourier mode.
The mass term $\lambda\chi$ and the source $\lambda\beta Y_{\rm qs}$
are algebraic (non-derivative) and do not contribute
to the principal symbol of the standalone $\chi$ equation when
$Y_{\rm qs}$ is treated as an external quasi-static background
functional.

For a Lorentzian metric $g_{\mu\nu}$ with signature
$(-,+,+,+)$, the symbol $\sigma(\xi) = 0$ defines
a cone with one timelike and three spacelike
directions.  This is the defining property of a
\emph{normally hyperbolic} operator.  By the standard
theory of linear wave equations on globally hyperbolic
spacetimes (see e.g.\ Wald, \textit{General Relativity}, 1984,
Ch.~10; Hawking \& Ellis, \textit{The Large Scale
Structure of Space-Time}, 1973, Ch.~7),
the Cauchy problem
\begin{equation}
\Box\chi - \lambda\chi = -\lambda\beta Y_{\rm qs}\,,\quad
\chi\big|_{\Sigma} = \chi_0\,,\quad
\dot\chi\big|_{\Sigma} = \chi_1\,,
\label{eq:cauchy}
\end{equation}
on any spacelike hypersurface $\Sigma$ admits a
unique, continuously dependent solution for smooth
initial data $(\chi_0,\chi_1)$ and source~$Y_{\rm qs}$.

\paragraph{Joint system $(\varphi,\chi)$.}
The preceding argument establishes the sourced KG problem for
$\chi$ on a prescribed quasi-static scalar background.  It does
\emph{not} by itself prove hyperbolicity of the fully varied
backreacting $(\varphi,\chi)$ system.  The reason is that
$Y_{\rm qs}$ depends on first derivatives of~$\varphi$; varying the
term $(\chi-\beta Y_{\rm qs})^2$ with respect to $\varphi$ can modify
the $\varphi$ principal block through derivative-coupling and kinetic
renormalization terms.  A direct two-field principal-symbol
calculation is therefore required before claiming full covariant
well-posedness of the backreacting system.  In the present appendix
we use Route~A only as a one-way quasi-static effective template.

\section{Normalization and parameter economy}\label{sec:norm}
\subsection{Closing the normalization via the cosmological MOND scale}
A key aim is to preserve the parameter economy of the main theory by fixing the free parameters through the MOND scale $a_0$, determined independently in the cosmological sector (the main theory) as
\begin{equation}
a_0 = \frac{cH_0}{2\pi}.
\label{eq:a0}
\end{equation}
The $2\pi$ factor follows from the standard Fourier convention:
\begin{equation}
k_H \equiv \frac{aH}{c},\qquad
\lambda_H \equiv \frac{2\pi a}{k_H}=\frac{2\pi c}{H},\qquad
\text{and we identify}\quad a_0 \equiv \frac{c^2}{\lambda_H}.
\end{equation}
The identification $a_0\equiv c^2/\lambda_H$ is a physical linkage (not a mathematical identity): it ties a characteristic cosmological wavelength scale to an acceleration scale in the emergent MOND sector.

\paragraph{Normalization target equation (explicit).}
To make the parameter-economy goal operational, we parametrize the halo-regime closure such that the asymptotic circular speed can be written as
\begin{equation}
V_{\rm flat}^4 \;=\; G M\,\mathcal{N}\!\left(\beta,\lambda\right)\,a_0,
\label{eq:normtarget}
\end{equation}
where $\mathcal{N}(\beta,\lambda)$ is the dimensionless normalization implied by the specific $\chi$ implementation.  In the covariant effective formulation of Sec.~\ref{sec:routeA_real}, the overdamped relaxation rate is Hubble-regulated ($\tau^{-1}=\lambda c^2/3H$), so in the static limit $\chi\to\chi_{\rm eq}=\beta Y_{\rm qs}$ the normalization depends only on~$\beta$; the stiffness~$\lambda$ controls only the \emph{approach} to equilibrium.  The \emph{parameter-free target} is $\mathcal{N}=1$, at which the BTFR normalisation would be inherited directly from $a_0$ with no additional fit parameter.  Whether the covariant closure actually realises $\mathcal{N}=1$---i.e.\ whether the correct $\beta$ is the one that already follows from $a_0=cH_0/2\pi$ with no residual multiplicative factor---is an open check and depends on the full bound-structure derivation in the companion MOND paper~\cite{Lotishvili2026MOND} and on the Route~B programme (Sec.~\ref{sec:ledger_app13}).  The non-circularity of the construction is nevertheless preserved: $a_0$ is determined first from the cosmological linkage in the main theory and only then used to set the galaxy-scale normalisation through Eq.~\eqref{eq:normtarget}; if the ultimate $\mathcal{N}$ differs from unity the discrepancy is absorbed into a single well-defined dimensionless closure ratio, not into an already-derived first-principles constant.

\section{Derived vs.\ Proposed: compact ledger}\label{sec:ledger_app13}
\begin{center}
\begin{tabular}{@{}p{0.27\linewidth}p{0.63\linewidth}@{}}
\toprule
\textbf{Derived (baseline)} &
Bi--conformal metric \eqref{eq:biconf_app13}; covariant action \eqref{eq:action_app13}; weak-field $\gamma_{\rm PPN}=1$; Shapiro 1PN with a 2PN discriminator; same-matter metric-sector linear-cosmology limit ($\alpha_i=0$ plus App.~9 scalar-stress decoupling; the main theory and CMB appendix); spectator sourcing \eqref{eq:deltaphi_app13}; MOND interpolating function $\mu(x)=x/(1+x)$ from the bound-structure Bernoulli$+$vortex closure (the companion MOND paper~\cite{Lotishvili2026MOND}).\\
\midrule
\textbf{Established (Route~A effective template)} &
Minimal augmented action~\eqref{eq:Schi} with $U=0$;
Klein--Gordon field equation~\eqref{eq:chi_KG};
ghost freedom (Sec.~\ref{sec:ghost_app13});
gradient stability and subluminal propagation
(Sec.~\ref{sec:gradient});
quasi-static reduction to first-order relaxation
with effective overdamped timescale
$\tau^{-1}=\lambda c^2/(3H)$~\eqref{eq:tau_derived};
normal hyperbolicity and well-posedness of the standalone sourced
$\chi$ equation for prescribed $Y_{\rm qs}$ (Sec.~\ref{sec:wellposed}).\\
\midrule
\textbf{Microscopic input (quantum extension~\cite{Lotishvili2026Q})} &
Microscopic tracking stiffness
$\lambda = \omega_0^2/c^2 = (m_0 c/\hbar)^2\,e^{\varphi}$
from the oscillon's fundamental quasi-normal mode
(Eq.~\ref{eq:lambda_micro}); this does not by itself derive the
macroscopic collective relaxation scale or the closure normalization
$\beta$.\\
\midrule
\textbf{Remaining open items} &
Route~B coarse-graining derivation of $\chi$ from
nonlinear $\varphi$ dynamics;
full backreacting $(\varphi,\chi)$ principal-symbol analysis;
macroscopic collective relaxation scale;
numerical halo profiles from the full
PDE~\eqref{eq:chi_flrw};
comparison of Hubble-regulated vs.\
gradient-dependent relaxation with observed
rotation curves;
determination of $\beta$ from first principles.\\
\bottomrule
\end{tabular}
\end{center}

\section{Falsifiable discriminants}\label{sec:pred}
The hysteresis closure is designed to be testable by comparing model variants. The most direct discriminants are:
\begin{itemize}
\item \textbf{Halo scaling and transition:} different $\tau(|\nabla\phi|)$ imply different radial slopes and transition radii.
\item \textbf{Lensing--dynamics relation:} because $\chi$ is tied to gradients of $\phi$, shear profiles and rotation curves are linked in a way that differs from NFW-based fits.
\item \textbf{Environment sensitivity:} the local-gradient dependence naturally introduces external-field sensitivity, providing a clear test with satellites and dwarfs.
\item \textbf{Redshift drift of normalization:} through $a_0(z)\propto H(z)$ (the main theory), the BTFR normalization is predicted to evolve with cosmic time.
\end{itemize}

\section{Notes for future work}\label{sec:notes}
The Route~A programme now provides a controlled effective template:
the $\chi$ field equation follows from the augmented action, and
ghost freedom, gradient stability, standalone KG hyperbolicity, and
the quasi-static overdamped reduction are established within the
one-way source treatment.  The remaining open items are:
\begin{enumerate}
\item \textbf{Route~B (coarse-graining):} derive $\chi$
  as an emergent slow variable from the nonlinear
  dynamics of~$\varphi$ around bound structures,
  connecting to the Bessel eigenvalue analysis of
  the companion MOND paper~\cite{Lotishvili2026MOND}.  (The quantum
  extension~\cite{Lotishvili2026Q} partially addresses
  the microscopic tracking scale through the oscillon
  quasi-normal mode structure,
  $\lambda_{\rm micro} = \omega_0^2/c^2$; a full
  coarse-graining from the classical nonlinear sector and the
  macroscopic collective relaxation scale remain open.)
\item \textbf{Backreacting principal symbol:} vary the full
  $(\varphi,\chi)$ action, including the derivative dependence of
  $Y_{\rm qs}$, and establish the conditions for covariant
  hyperbolicity of the coupled system.
\item \textbf{Numerical profiles:} solve the full
  PDE~\eqref{eq:chi_flrw} for realistic galaxy
  potentials and compare the resulting rotation curves
  and lensing profiles with observations.
\item \textbf{Relaxation-rate discriminant:} the
  Hubble-regulated rate
  $\tau^{-1} = \lambda c^2/(3H)$ and the original
  gradient-dependent rate
  $\tau^{-1} = \gamma|\nabla\varphi|$ predict
  different outer-halo slopes; quantify this difference
  for observational tests.
\end{enumerate}

\section*{Takeaway}
The $\chi$ hysteresis sector is now placed on a
covariant effective-template footing.  Starting from the minimal
augmented action~\eqref{eq:Schi} with $U = 0$ and prescribed
quasi-static source $Y_{\rm qs}$:
(i)~the field equation~\eqref{eq:chi_KG} is a standard
massive Klein--Gordon equation sourced by quasi-static $\varphi$
gradients in the bound-structure sector;
(ii)~the sector is ghost-free, gradient-stable,
and normally hyperbolic as a standalone sourced KG problem;
(iii)~in the quasi-static, overdamped limit the
covariant equation reduces to the phenomenological
relaxation law with an effective Hubble-friction timescale
$\tau^{-1} = \lambda c^2/(3H)$;
(iv)~the microscopic tracking stiffness
$\lambda_{\rm micro} = \omega_0^2/c^2$ is tied to the oscillon
quasi-normal mode structure in the quantum
extension~\cite{Lotishvili2026Q}
(Eq.~\ref{eq:lambda_micro}), while the macroscopic collective
relaxation scale and the closure normalization $\beta$ remain open
Route~B/coarse-graining tasks.  The BTFR normalization target is
anchored by the independently determined $a_0 = cH/(2\pi)$, but its
full realization remains a bound-structure closure check.
\end{document}
